\section{Predict, run, explain with the prompts}
\begin{frame}{Prompt 01a: model the graph}
\princ{MODEL VS PROCEDURE}
\promptbox{01}{a}
\predict{Before running: which exact seven edges should the test require?}
\see{the model is written down before the code is trusted.}
\note{Give about three minutes. Students predict the seven vertices and edges, then run prompt-a. Listen for the phrase ``A must be completed before B.'' The model comes before the procedure.}
\end{frame}

\begin{frame}{Prompt 01b: test the graph}
\princ{MODEL VS PROCEDURE}
\promptbox{01}{b}
\see{the exact-set test checks the graph without allowing later code to change it.}
\note{Students run prompt-b after prompt-a. The exact-set test is the visible check; do not allow later code to silently change the data.}
\end{frame}

\begin{frame}{Prompt 02a: draw the directory tree}
\princ{TREE STRUCTURE}
\promptbox{02}{a}
\predict{Before running: which edge changes meaning when it becomes a parent-child line?}
\see{the picture is a containment view, so its edges need a new definition.}
\note{Give about four minutes. Have students convert the graph into the fixed directory-style tree and name the edge meaning before they judge the drawing.}
\end{frame}

\begin{frame}{Prompt 02b: test the tree}
\princ{TREE STRUCTURE}
\promptbox{02}{b}
\mistake{A neat directory picture is not automatically a complete dependency graph.}
\note{Students run the structural checks: root, one parent, seven nodes, and acyclicity. The test checks the tree shape, not whether it preserves every dependency.}
\end{frame}

\begin{frame}{Prompt 03a: trace the tree}
\princ{TRAVERSAL}
\promptbox{03}{a}
\predict{Before running: which list should put build before requirements?}
\see{the three traversal rules answer three different visit-order questions.}
\note{Give about four minutes. Students implement preorder, postorder, and breadth-first traversal. Remind them that the fixed child order is part of the input.}
\end{frame}

\begin{frame}{Prompt 03b: test the traces}
\princ{TRAVERSAL}
\promptbox{03}{b}
\see{the breadth-first list is the level-by-level list generated from the fixed order.}
\note{Students compare the exact assertions. If breadth-first looks like preorder, ask where the queue is. Inorder is not part of this general rooted-tree lesson.}
\end{frame}

\begin{frame}{Prompt 04a: trace a decision tree}
\princ{DECISION RULE}
\promptbox{04}{a}
\predict{Before running: should yes or no appear in the node trace?}
\see{the route records nodes; yes and no label which edge leads to the next node.}
\note{Give about four minutes. Students model the eligible-plus-aid route, keeping yes/no as edge labels. No real intake data is needed or allowed.}
\end{frame}

\begin{frame}{Prompt 04b: test the decision tree}
\princ{DECISION RULE}
\promptbox{04}{b}
\see{the trace ends at aid route; the review-before-use condition is a separate note.}
\note{Students test the three decision traces and identify what the tree does not show: the language-access review must happen before the aid route is usable.}
\end{frame}

\begin{frame}{Prompt 05a: audit the cross-link}
\princ{CHECK THE WORK}
\promptbox{05}{a}
\predict{Before running: what dropped edge set do you predict?}
\see{set subtraction compares relationships, not just the names of nodes.}
\note{Give about four minutes. Students compare graph edges with tree edges and predict the one missing relationship before running the audit.}
\end{frame}

\begin{frame}{Prompt 05b: add the audit test}
\princ{CHECK THE WORK}
\promptbox{05}{b}
\see{the audit names the actual dropped relation; the strict expected failure stays visible.}
\note{Students add the regression test and keep the strict expected failure. Explain that an expected failure is evidence about the false claim, not a broken suite.}
\end{frame}

\begin{frame}{Prompt 06a: check a search claim}
\princ{HIERARCHY LIMIT}
\promptbox{06}{a}
\predict{Before running: can finding every known node prove every dependency is present?}
\see{a node search asks whether an item is present; an edge audit asks whether a relationship is present.}
\note{Give about four minutes. Students search the seven synthetic tree nodes and include the missing language-access review case.}
\end{frame}

\begin{frame}{Prompt 06b: test the stronger claim}
\princ{HIERARCHY LIMIT}
\promptbox{06}{b}
\mistake{A successful node search is not a relationship audit.}
\note{Students compare node coverage with the dropped-edge audit. The intended conclusion is no: finding nodes does not prove that all dependency edges were preserved.}
\end{frame}

\begin{frame}{Prompt 07a: explain the limit}
\princ{CHECK THE WORK}
\promptbox{07}{a}
\predict{Before running: what evidence will support the one-sentence conclusion?}
\see{your conclusion matches the traces, the dropped dependency, and the limits of those checks.}
\note{Give about five minutes. Students run the full suite and print both project and decision traces plus the dropped set. Require a student-written sentence and have them label direct checks versus model limitations.}
\end{frame}
